% !TeX spellcheck = it_IT

\documentclass[aspectratio=169, handout]{beamer}

\usepackage[utf8]{inputenc}
\usepackage[italian]{babel}
\usepackage{multicol}

\usetheme[titleformat=smallcaps, numbering=fraction, background=light, progressbar=frametitle]{metropolis}

\usepackage{mathtools}
\DeclarePairedDelimiter{\ceil}{\lceil}{\rceil}

\title{Esercitazioni di Informatica B}
\subtitle{Introduzione a Matlab}

\author{Stefano Cereda\\
	stefano.cereda@polimi.it
}
\date{12/11/2019}
\institute[PoliMi]{Politecnico Milano}
\logo{\includegraphics[width=10mm]{../logopolimi}}

\setbeamercovered{invisible}

\makeindex

\begin{document}
	\begin{frame}
	\maketitle
\end{frame}

%\begin{frame}<handout:0>{Conclusione esercizi struct}
%Completare gli esercizi sulle struct in C (anagrafe)
%\end{frame}


\begin{frame}{Matlab/Octave}
    \begin{itemize}
        \item Installate Matlab il prima possibile: \url{https://www.software.polimi.it/}
        \item In alternativa, potete utilizzare Octave: \url{https://www.gnu.org/software/octave/}
        \item I due software hanno interfacce simili, con esplora directory, editor di testo, finestra comandi e
            workspace.
        \item Matlab contiene funzionalità aggiuntive che userete in altri corsi.
    \end{itemize}
\end{frame}

\begin{frame}{Costrutto if: equazioni di secondo grado}
Si chieda l’inserimento dei tre coefficienti di una equazione di secondo grado: $ax^2+bx+c=0$ e si stampi il valore delle radici reali.

%Script: commenti(nome e descrizione; data e autori; def variabili), input, calcolo, output.

\pause
Cosa succede se chiediamo di trovare radici complesse?

\pause
Cosa cambia fra scrivere uno script ed inserire direttamente i comandi?

\pause
Come funziona il debugger di Matlab?
\end{frame}

\begin{frame}{Costrutto while: fattoriale}
Si scriva uno script che chieda l'inserimento di un numero $n$ e ne calcoli e visualizzi il fattoriale utilizzando un ciclo while.
\end{frame}

\begin{frame}{Matrice triangolare}
Scrivere uno script che svolga le seguenti operazioni:
\begin{enumerate}
	\item Definire una matrice quadrata di interi di dimensione $N\times N$, con $N=5$.
	\item Acquisire da tastiera una sequenza di valori interi da memorizzare nella porzione triangolare superiore della
        matrice. La parte triangolare superiore della matrice è costituita dalle celle poste sulla diagonale principale
        e al di sopra di tale diagonale.
	\item Copiare in un array di dimensioni opportune tutti i valori della porzione triangolare di matrice che sono maggiori di -10.
	\item Stampare il contenuto dell'array.
\end{enumerate}
\end{frame}

\begin{frame}{Grafici: equazioni di secondo grado}
Si chieda l’inserimento dei tre coefficienti di una equazione di secondo grado: $ax^2+bx+c=0$ e si dia una rappresentazione della curva.
\end{frame}

\begin{frame}{Rotazione di matrice}
Sc scriva una script in Matlab/Octave che data una matrice $a$ di dimensioni $N\times N$ crea una nuova matrice $b$ ruotata di 90 gradi in senso antiorario rispetto ad $a$.

Si consideri $N=4$ e la matrice $a$ inizializzata con i valori:

\centering
\begin{tabular}{|cccc|}
	\hline
	1&2&3&4\\
	2&3&4&5\\
	6&7&8&9\\
	0&0&0&0\\
	\hline
\end{tabular}
\end{frame}
\end{document}
